\chapter{Conclusion}
\label{chap:conclusion}

This project successfully demonstrated the complete RTL-to-GDSII implementation of a UART (Universal Asynchronous Receiver-Transmitter) IP core using the open-source OpenLane2 ASIC design flow and the SkyWater SKY130 Process Design Kit (PDK). The UART was designed in Verilog HDL using a modular architecture consisting of a Baud Rate Generator, UART Transmitter, UART Receiver, and a top-level integration module. Functional verification through simulation confirmed the correct operation of each module as well as the complete UART system.

The verified RTL was synthesized using Yosys and subsequently implemented through the OpenROAD physical design flow. The implementation included floorplanning, standard-cell placement, clock tree synthesis, global routing, and detailed routing, resulting in a complete physical layout of the UART IP core. Throughout the implementation process, the design achieved successful placement, balanced clock distribution, congestion-free routing, and efficient utilization of the available routing resources.

Post-layout verification was performed using OpenRCX, OpenSTA, Magic, and Netgen. RC extraction, timing analysis, IR-drop analysis, Design Rule Checking (DRC), Layout Versus Schematic (LVS) verification, and manufacturability assessment were all completed successfully. The generated layout reported zero routing overflow, zero antenna violations, zero DRC violations, successful LVS matching, and negligible IR drop, demonstrating the correctness and fabrication readiness of the design.

Overall, this project provides practical experience with the complete digital ASIC implementation flow using entirely open-source EDA tools. It demonstrates the transformation of a synthesizable RTL design into a verified GDSII layout while exposing every major stage of modern ASIC design, including synthesis, physical implementation, verification, and signoff.

\section*{Future Work}

The current UART IP core can be extended to support additional communication features and more complex System-on-Chip (SoC) designs. Possible future enhancements include:

\begin{itemize}
    \item Configurable baud rates and programmable data formats.
    \item Support for parity generation and parity error detection.
    \item FIFO-based transmit and receive buffers for improved throughput.
    \item Interrupt-driven UART operation for processor-based systems.
    \item Integration with standard on-chip buses such as APB or AXI-Lite.
    \item FPGA prototyping followed by fabrication through an MPW shuttle run.
\end{itemize}